\scanpage{289}
\mspage{15}
% archive sheet number: 288

surjectif, donc $\exists\,s$ tel que $u_2s$ soit homotope à $u_1$ dans $\mathrm{Cat}$, ce qui implique $u_2s=u_1$ dans $\mathrm{Hot}_{\unsure{W|w}}$, \unsure{qed|qed.}.

\uline{Commentaire final}

L'axiome de connexité\inserted{Loc 5} diffère de tous les autres axiomes Loc 1 à Loc 8 (\unsure{mais|même} à l'exception de Loc 8 bis d'accessibilité) par le fait que ce n'est pas un axiome de stabilité.

Les autres axiomes Loc 1 à Loc 8 (et aussi Loc 8 bis d'accessibilité) et -- les axiomes \unsure{ultérieurs} \unsure{à considérer|de continuité}$^{\ast}$, \unsure{sont} satisfaits si $W=\unsure{\mathrm{Fl}|\mathrm{El}}\,M$ -- l'axiome de connexité est donc le seul qui s'oppose à la \uline{trivialité} de la théorie homotopique définie par $W$. C'est dire qu'il ne peut être conséquence de l'\unsure{ens.|un} des autres.
% slanted
\margin{il n'y en a en fait qu'un, grâce à la ``\unsure{providentielle}'' \unsure{thèse|théorie} de K.\unsure{S.} Brown !}

Par contre, l'axiome $\mathrm{Loc}(4)$ des sommes directes n'est pas indépendant des autres. Signalons à ce propos (rappel)
\[
\left\{\begin{aligned}
&\mathrm{Loc}(4a_0)\ \text{est conséquence de l'axiome $\mathrm{Loc}(6a)$}\\
&\hspace{3em}\text{des carrés cocartésiens (forme \unsure{directe|ouverte})}\\
&\mathrm{Loc}(4\unsure{a})\ \text{est conséquence de $\mathrm{Loc}(6a)$ précédent,}\\
&\hspace{3em}\text{$+$ l'axiome des limites inductives}\\
&\hspace{3em}\text{filtrantes $\mathrm{Loc}(8)$}\\
&\mathrm{Loc}(4b)\ \text{est conséquence de l'axiome $\mathrm{Loc}(7\unsure{)}$}\\
&\hspace{3em}\text{des fibrations $+$ \deleted{\unsure{donc}}}\\
&\mathrm{Loc}(4)\ \text{est conséquence de $\mathrm{Loc}(6a,7,8)$.}\\
&\hspace{3em}\text{(un point pour une vache,}\\
&\hspace{6em}\text{il est vrai\ldots)}
\end{aligned}\right.
\]
